\chapter{Qualité et intégration continue}
\label{chap-qualite}

Le cours pose une règle de méthode : une fonction écrite est une fonction testée. La
version~1 ne pouvait pas l'appliquer, puisque ses règles ne s'exécutaient pas sans fenêtre.
La version~2 en fait une chaîne automatique, exécutée à chaque modification du code.

\section{Stratégie de test}

\begin{table}[ht]
    \centering
    \small
    \begin{tabularx}{\textwidth}{@{}l>{\raggedright\arraybackslash}X@{}}
        \toprule
        \textbf{Niveau}     & \textbf{Ce qui est vérifié, et comment}                                  \\
        \midrule
        Unitaire            & 54 cas répartis en cinq programmes : chaque fonction publique du moteur,
        chaque chemin d'erreur.                                                                        \\
        \addlinespace
        Propriétés          & \num{2250} parties complètes jouées par le moteur contre lui-même ; tous
        les invariants sont recalculés \emph{après chaque tir}.                                        \\
        \addlinespace
        Stratégies          & Les trois niveaux doivent rester ordonnés par force mesurée, avec une
        marge ; un bateau touché doit être poursuivi.                                                  \\
        \addlinespace
        Reproductibilité    & Une simulation à graine fixée est comparée, octet pour octet, à un
        fichier de référence versionné.                                                                \\
        \addlinespace
        Interface terminal  & Une session scriptée est comparée à la transcription attendue.           \\
        \addlinespace
        Interface graphique & Une partie entière est jouée à travers les vrais boutons GTK,
        tout avertissement de GTK étant fatal.                                                         \\
        \bottomrule
    \end{tabularx}
    \caption{Les niveaux de test.}
    \label{tab-tests}
\end{table}

\subsection{Un harnais de test de cent lignes}

Le moteur ne dépend de rien ; il aurait été dommage que ses tests imposent une bibliothèque.
Le harnais \fichier{tests/bs\_test.h} tient en une centaine de lignes : des macros
\code{TEST}, \code{CHECK}, \code{CHECK\_EQ} et \code{REQUIRE}, une table de cas et une
fonction principale qui produit une sortie au format TAP (\cref{lst-test}).

\begin{lstlisting}[basicstyle=\ttfamily\footnotesize,
                   caption={Un test unitaire : un tir refusé ne modifie rien (exigence~C4).},
                   label={lst-test}]
TEST(refused_shots_leave_the_game_untouched)
{
    bs_game game;
    bs_coord deck;
    REQUIRE(bs_game_init(&game, BS_DIFFICULTY_EASY, 2) == BS_OK);
    REQUIRE(find_target(&game, true, &deck));
    REQUIRE(bs_game_human_fire(&game, deck, NULL) == BS_OK);

    CHECK_EQ(BS_ERR_ALREADY_SHOT, bs_game_human_fire(&game, deck, NULL));
    CHECK_EQ(BS_ERR_OUT_OF_BOUNDS, bs_game_human_fire(&game, bs_coord_make(10, 0), NULL));
    CHECK_EQ(1, bs_game_stats(&game, BS_PLAYER_HUMAN)->shots);
    CHECK_EQ(BS_PHASE_HUMAN_TURN, bs_game_phase(&game));
}
\end{lstlisting}

\subsection{Tests de propriétés}
\label{sec-proprietes}

Un test unitaire vérifie un exemple choisi par son auteur. Un test de propriété vérifie une
\emph{loi} sur des milliers d'exemples que personne n'a choisis. Le programme
\fichier{test\_simulation} joue \num{2250} parties --- 250 pour chacune des neuf
combinaisons de niveaux --- et, après chaque tir, recalcule tout ce que le moteur tient à
jour pour le comparer à ce qu'il annonce :

\begin{itemize}
    \item le nombre de ponts et de bateaux à flot, recompté case par case ;
    \item la cohérence des états : aucune case d'eau n'est annoncée touchée, aucun bateau
          n'est annoncé dans l'eau, une case est \code{SUNK} exactement quand son bateau est
          coulé ;
    \item le score, qui doit se lire intégralement sur le plateau ;
    \item le tour de jeu : la main ne change que sur un tir manqué ;
    \item la fin : jamais deux flottes coulées, un vainqueur correctement désigné.
\end{itemize}

Exécuté sous les \emph{sanitizers} (\cref{sec-sanitizers}), ce test fait aussi office de
campagne de \emph{fuzzing} sur l'ensemble du moteur.

\subsection{Tester l'interface graphique pour de vrai}

L'option \code{-{}-autoplay} fait jouer l'ordinateur des deux côtés. Elle aurait pu envoyer
directement les coups du \enquote{joueur} au modèle ; elle n'aurait alors rien testé de
l'interface. À la place, le contrôleur demande à la vue de \emph{cliquer sur le bouton} de
la case choisie. Chaque coup parcourt ainsi le trajet complet d'une action de
l'utilisateur :

\begin{center}
    bouton GTK $\;\rightarrow\;$ signal \enquote{clicked} $\;\rightarrow\;$ vue
    $\;\rightarrow\;$ contrôleur $\;\rightarrow\;$ modèle $\;\rightarrow\;$ synchronisation.
\end{center}

En fin de partie, le test répond \enquote{quitter} à la boîte de dialogue, là encore comme
le ferait un clic. Dans la chaîne d'intégration continue, ce test tourne sous un serveur
d'affichage virtuel (Xvfb), avec la variable \code{G\_DEBUG=fatal-criticals} qui transforme
toute anomalie signalée par GTK en échec. Les effets sonores y sont réellement décodés, vers
une sortie factice, et un son qui ne peut pas être joué fait échouer le test. Celui-ci
enregistre enfin une capture de la fenêtre finale comme pièce jointe de la construction. L'animation de la page d'accueil du dépôt est
produite de cette manière, et non filmée à la main.

\subsection{Reproductibilité entre plates-formes}

Le moteur n'utilise que de l'arithmétique entière. Une simulation à graine fixée doit donc
donner le même histogramme avec tous les compilateurs, sur tous les systèmes et toutes les
architectures. Un test compare la sortie du simulateur à un fichier versionné
(\fichier{tests/data/simulate\_seed2023.csv}). Il échoue si le générateur aléatoire, le
placement des flottes ou une stratégie change de comportement --- que ce soit par erreur, ou
volontairement, auquel cas le fichier de référence doit être régénéré et le changement
devient visible dans l'historique (exigence~C2).

\section{Analyse du code}

\subsection{Avertissements}

Le code est compilé avec un jeu d'avertissements étendu (\code{-Wall -Wextra -Wpedantic
    -Wshadow -Wconversion -Wsign-conversion -Wstrict-prototypes}\dots) et tout avertissement est
une erreur (exigence~C5). La version~1 ne demandait que \code{-Wall}.

\subsection{Analyse dynamique}
\label{sec-sanitizers}

Toute la suite de tests est exécutée une seconde fois avec
AddressSanitizer~\cite{serebryany2012} et UndefinedBehaviorSanitizer, qui instrumentent le
programme pour détecter à l'exécution les accès mémoire invalides et les comportements
indéfinis. La lecture hors limites de la version~1 (\cref{chap-adversaire}) est exactement
le type de défaut que cet outil révèle au premier lancement.

\subsection{Analyse statique et couverture}

Trois analyseurs examinent le code sans l'exécuter : \emph{clang-tidy}, \emph{cppcheck} et
\emph{CodeQL}. La mise en forme est vérifiée par \emph{clang-format}. Enfin, la couverture
des tests est mesurée sur le moteur (\cref{tab-couverture}) ; la construction échoue si
elle passe sous \SI{95}{\percent}.

\begin{table}[ht]
    \centering
    \begin{tabular}{@{}l S[table-format=3.0] S[table-format=3.2] l@{}}
        \toprule
        \textbf{Fichier}   & {\textbf{Lignes}} & {\textbf{Couvertes (\si{\percent})}} &
        \textbf{Lignes non couvertes}                                                                               \\
        \midrule
        \fichier{ai.c}     & 130               & 100.00                               &                             \\
        \fichier{board.c}  & 186               & 97.85                                & impasse de placement        \\
        \fichier{common.c} & 14                & 100.00                               &                             \\
        \fichier{game.c}   & 101               & 98.02                                & échec du placement, propagé \\
        \fichier{rng.c}    & 40                & 100.00                               &                             \\
        \midrule
        \textbf{Moteur}    & {\bfseries 471}   & {\bfseries 98,73}                    &                             \\
        \bottomrule
    \end{tabular}
    \caption{Couverture de lignes du moteur par ses tests. Les six lignes non couvertes
        traitent toutes le cas, jamais observé, où une flotte ne peut pas être placée.}
    \label{tab-couverture}
\end{table}

\section{Intégration continue}
\label{sec-ci}

\subsection{Une seule description de la construction}

Le projet est construit par CMake. Chaque configuration utile --- développement, version
optimisée, \emph{sanitizers}, couverture, construction sans interface graphique --- est un
\emph{préréglage} nommé du fichier \fichier{CMakePresets.json}~\cite{cmakepresets}. La chaîne d'intégration
continue n'exécute rien d'autre que

\begin{center}
    \code{cmake -{}-preset X} \quad\code{cmake -{}-build -{}-preset X}\quad
    \code{ctest -{}-preset X}
\end{center}

et aucune option de compilation n'est écrite dans un fichier de la chaîne. Un échec sur le
serveur se reproduit donc sur un poste de travail par la même commande.

\subsection{La chaîne}

\begin{figure}[ht]
    \centering
    \begin{tikzpicture}[node distance=4mm and 7mm,
            tache/.style={boite, font=\sffamily\footnotesize, minimum width=3.5cm,
                    minimum height=8.5mm, inner sep=3pt},
            titre/.style={font=\sffamily\small\bfseries, text=marine}]
        \node[coeur, minimum width=2.0cm, minimum height=1.1cm] (push) {\texttt{git push}};

        \node[tache, right=13mm of push, yshift=24mm] (fmt) {format\\[-1pt]\scriptsize clang-format};
        \node[tache, below=of fmt] (build) {build ($\times 5$)\\[-1pt]\scriptsize GCC, Clang, Apple Clang};
        \node[tache, below=of build] (head) {headless\\[-1pt]\scriptsize Debian, sans GTK};
        \node[tache, right=of fmt] (asan) {sanitizers\\[-1pt]\scriptsize ASan + UBSan};
        \node[tache, below=of asan] (cov) {coverage\\[-1pt]\scriptsize seuil \SI{95}{\percent}};
        \node[tache, below=of cov] (lint) {lint\\[-1pt]\scriptsize clang-tidy, cppcheck};
        \node[tache, below=of head, fill=brume] (codeql) {CodeQL};
        \node[tache, below=of lint, fill=brume] (report) {rapport\\[-1pt]\scriptsize \LaTeX{} $\rightarrow$ PDF};

        \begin{scope}[on background layer]
            \node[draw=ocean, dashed, rounded corners=4pt, inner sep=6pt, fill=ecume!50,
                fit=(fmt)(lint)(build)(head)(asan)(cov)] (ci) {};
        \end{scope}
        \node[titre, above=0.5mm of ci.north] {Intégration continue (\fichier{ci.yml})};

        \draw[fleche] (push.east) -- ++(5mm,0) |- (ci.west);
        \draw[fleche] (push.east) -- ++(5mm,0) |- (codeql.west);

        \node[coeur, below=26mm of push, minimum width=2.0cm, minimum height=1.1cm] (tag)
        {étiquette\\[-1pt]\texttt{v*}};
        \node[tache, minimum width=7.7cm, fill=sable, draw=corail] (rel)
        at ($(codeql.south)!0.5!(report.south) + (0,-13mm)$)
        {livraison (\fichier{release.yml}) : tests, paquets \texttt{.tar.gz} et \texttt{.deb},\\[-1pt]
            sommes de contrôle, attestation de provenance, PDF du rapport};
        \draw[fleche, draw=corail] (tag.east) -- ++(5mm,0) |- (rel.west);
    \end{tikzpicture}
    \caption{La chaîne d'intégration et de livraison continues (GitHub Actions).}
    \label{fig-ci}
\end{figure}

À chaque \emph{push} et à chaque demande de fusion, six tâches s'exécutent en parallèle
(\cref{fig-ci}) :

\begin{description}[style=nextline, leftmargin=1.2em]
    \item[format] Vérifie la mise en forme de tous les fichiers source.
    \item[build] Compile et teste avec GCC et Clang sous Ubuntu, en mode \emph{Debug} et
          \emph{Release}, et avec Apple Clang sous macOS ; sous Linux, joue en plus la partie de
          bout en bout dans l'interface graphique et conserve la capture d'écran.
    \item[headless] Compile et teste le moteur et l'interface terminal dans un conteneur
          Debian où GTK n'est pas installé : c'est la preuve de l'exigence~C1. Ce conteneur
          fournit aussi la plus ancienne version de CMake prise en charge.
    \item[sanitizers] Rejoue tous les tests, interface graphique comprise, sous ASan et
          UBSan.
    \item[coverage] Mesure la couverture du moteur et applique le seuil.
    \item[lint] Lance clang-tidy et cppcheck.
\end{description}

Deux autres chaînes complètent le dispositif : l'analyse de sécurité CodeQL, et la
compilation du présent rapport, dont le PDF est publié comme pièce jointe.

\subsection{Livraison}

Poser une étiquette de version \code{v2.0.0} déclenche la livraison. La chaîne recompile en
mode optimisé, rejoue les tests, produit les paquets (\code{.tar.gz}, \code{.deb}), calcule
leurs sommes de contrôle SHA-256, signe une \emph{attestation de provenance} qui relie
chaque fichier publié au \emph{commit} et à l'exécution qui l'ont produit, et publie le
tout avec le PDF du rapport et les notes de version extraites du journal des modifications.

\subsection{Chaîne d'approvisionnement}

Les actions tierces utilisées par les chaînes sont épinglées par l'empreinte de leur
\emph{commit} et non par un numéro de version, qu'un tiers pourrait déplacer ; un robot
propose chaque mois leurs mises à jour. Chaque tâche ne reçoit que les permissions dont elle
a besoin, la lecture seule par défaut.

\section{Conventions}

\begin{itemize}
    \item Les messages de \emph{commit} suivent les \emph{Conventional Commits}
          (\code{feat(ai): \dots}, \code{fix(ui): \dots}).
    \item Les versions suivent le versionnage sémantique~\cite{semver} ; le fichier
          \fichier{CHANGELOG.md} décrit chaque version au format \emph{Keep a Changelog}.
    \item Les choix d'architecture sont consignés dans de courtes fiches de décision
          (\fichier{docs/adr/}) : le contexte, la décision, ses conséquences.
\end{itemize}
